% Einheit 1 von 3 – Potenzen (bank/potenzen-wurzeln/e1.jsonl, zusammenbau v0.1)
%% TODO Zweigzeile Teil 1: was hier gelernt wird (Satz in Schülersprache aus dem Einheitstitel) – Einheit 1
\einheitenkopf[e1]{Einheit 1 von 3 · Potenzen}
\zweigzeile{ab Kl. 5, je nach Buch bis Kl. 9 · P10 oft}
\setcounter{aufgabe}{11}

%% TODO Titel: Ich-kann-Satz fehlt in der Bank (Platzhalter: Kettenname) – Nr. 12
%% TODO Anweisung: ein Satz über der ersten Teilaufgabe fehlt in der Bank – Nr. 12
\begin{aufgabe}{Plus oder minus?}
\begin{teile}
\teil $(-5)^4$ – plus oder minus?\\ \kreuz{positiv} \\ \kreuz{negativ}
\end{teile}
\end{aufgabe}

%% TODO Titel: Ich-kann-Satz fehlt in der Bank (Platzhalter: Kettenname) – Nr. 13
%% TODO Anweisung: ein Satz über der ersten Teilaufgabe fehlt in der Bank – Nr. 13
\begin{aufgabe}{Bruch oder minus?}
\begin{teile}
\teil $7^{-2}$ – Bruch oder negative Zahl?\\ \kreuz{Bruch} \\ \kreuz{negative Zahl}
\end{teile}
\end{aufgabe}

%% TODO Titel: Ich-kann-Satz fehlt in der Bank (Platzhalter: Kettenname) – Nr. 14
%% TODO Anweisung: ein Satz über der ersten Teilaufgabe fehlt in der Bank – Nr. 14
\begin{aufgabe}{Potenzen}
\begin{teile}
\teil $7^4$ – hoch oder mal?\\ \kreuz{hoch} \\ \kreuz{mal}
\end{teile}
%% TODO Darstellung für schwach fehlt in der Bank (potenzen-wurzeln-e1-k3-s1-v1; Abschnitt „Für schwache Schüler“)
\swz{$2^6$ als Malkette – wie viel?}{}
%% TODO Darstellung für schwach fehlt in der Bank (potenzen-wurzeln-e1-k3-s1-v2; Abschnitt „Für schwache Schüler“)
\swz{$7^3$ als Malkette – wie viel?}{}
%% TODO Darstellung für schwach fehlt in der Bank (potenzen-wurzeln-e1-k3-s1-v3; Abschnitt „Für schwache Schüler“)
\swz{$6 \cdot 6 \cdot 6$ als Potenz – wie viel?}{}
%% TODO Darstellung für schwach fehlt in der Bank (potenzen-wurzeln-e1-k3-s1-v4; Abschnitt „Für schwache Schüler“)
\swz{$11 \cdot 11$ als Potenz – wie viel?}{}
%% TODO Darstellung für schwach fehlt in der Bank (potenzen-wurzeln-e1-k3-s1-v5; Abschnitt „Für schwache Schüler“)
\swz{$3^5$ als Malkette – wie viel?}{}
\swfrage{Was bleibt gleich, was ändert sich?}
%% TODO Darstellung für schwach fehlt in der Bank (potenzen-wurzeln-e1-k3-s2-v1; Abschnitt „Für schwache Schüler“)
\swz{$7^2$ und $7 \cdot 2$ – wie viel jeweils?}{}
%% TODO Darstellung für schwach fehlt in der Bank (potenzen-wurzeln-e1-k3-s3-v1; Abschnitt „Für schwache Schüler“)
\swz{$14^2$ – wie viel?}{}
%% TODO Darstellung für schwach fehlt in der Bank (potenzen-wurzeln-e1-k3-s4-v1; Abschnitt „Für schwache Schüler“)
\swz{$23^2$ – wie viel? Rechne mit dem Taschenrechner.}{}
%% TODO Darstellung für schwach fehlt in der Bank (potenzen-wurzeln-e1-k3-s5-v1; Abschnitt „Für schwache Schüler“)
\swz{$(-4)^3$ – wie viel?}{}
%% TODO Darstellung für schwach fehlt in der Bank (potenzen-wurzeln-e1-k3-s6-v1; Abschnitt „Für schwache Schüler“)
\swz{$-8^2$ – wie viel?}{}
%% TODO Darstellung für schwach fehlt in der Bank (potenzen-wurzeln-e1-k3-s7-v1; Abschnitt „Für schwache Schüler“)
\swz{$0{,}6^2$ – wie viel?}{}
%% TODO Darstellung für schwach fehlt in der Bank (potenzen-wurzeln-e1-k3-s8-v1; Abschnitt „Für schwache Schüler“)
\swz{$2^7$ oder $5^3$ – welche Zahl ist größer?}{}
%% TODO Darstellung für schwach fehlt in der Bank (potenzen-wurzeln-e1-k3-s9-v1; Abschnitt „Für schwache Schüler“)
\swz{$5^3$, $2^6$, $9^2$ – welche Zahl ist die größte? Unterstreiche sie.}{}
%% TODO Darstellung für schwach fehlt in der Bank (potenzen-wurzeln-e1-k3-s10-v1; Abschnitt „Für schwache Schüler“)
\swz{$5^x = 625$ – welche Zahl ist $x$?}{}
%% TODO Darstellung für schwach fehlt in der Bank (potenzen-wurzeln-e1-k3-s11-v1; Abschnitt „Für schwache Schüler“)
\swz{$4^x = 4\,096$ – welche Zahl ist $x$?}{}
%% TODO Darstellung für schwach fehlt in der Bank (potenzen-wurzeln-e1-k3-s12-v1; Abschnitt „Für schwache Schüler“)
\swz{$10^x = 10\,000$ – welche Zahl ist $x$?}{}
%% TODO Darstellung für schwach fehlt in der Bank (potenzen-wurzeln-e1-k3-s13-v1; Abschnitt „Für schwache Schüler“)
\swz{$4^{-2}$ – als Bruch und als Dezimalzahl?}{}
%% TODO Darstellung für schwach fehlt in der Bank (potenzen-wurzeln-e1-k3-s14-v1; Abschnitt „Für schwache Schüler“)
\swz{$4^{-3} \;\square\; 0{,}02$ – $<$, $=$ oder $>$?}{}
\swz[3]{$2^x = 512$ – welche Zahl ist $x$? \hfill (P10 2020 OS)}{}
\end{aufgabe}

%% TODO Titel: Ich-kann-Satz fehlt in der Bank (Platzhalter: Kettenname) – Nr. 15
%% TODO Anweisung: ein Satz über der ersten Teilaufgabe fehlt in der Bank – Nr. 15
\begin{aufgabe}{Potenzen von Brüchen}
%% TODO Darstellung für schwach fehlt in der Bank (potenzen-wurzeln-e1-k4-s1-v1; Abschnitt „Für schwache Schüler“)
\swz{$\left(\frac{2}{5}\right)^2$ – wie viel?}{}
\end{aufgabe}

%% TODO Titel: Ich-kann-Satz fehlt in der Bank (Platzhalter: Kettenname) – Nr. 16
%% TODO Anweisung: ein Satz über der ersten Teilaufgabe fehlt in der Bank – Nr. 16
\begin{aufgabe}{Potenz gegen Dezimalzahl und Prozent vergleichen}
%% TODO Darstellung für schwach fehlt in der Bank (potenzen-wurzeln-e1-k5-s1-v1; Abschnitt „Für schwache Schüler“)
\swz{$0{,}5^2 \;\square\; 20\,\%$ – $<$, $=$ oder $>$?}{}
\end{aufgabe}

%% TODO Titel: Ich-kann-Satz fehlt in der Bank (Platzhalter: Kettenname) – Nr. 17
%% TODO Anweisung: ein Satz über der ersten Teilaufgabe fehlt in der Bank – Nr. 17
\begin{aufgabe}{hoch null}
%% TODO Darstellung für schwach fehlt in der Bank (potenzen-wurzeln-e1-k6-s1-v1; Abschnitt „Für schwache Schüler“)
\swz{$8^0$ – wie viel?}{}
\end{aufgabe}

%% TODO Titel: Ich-kann-Satz fehlt in der Bank (Platzhalter: Kettenname) – Nr. 18
%% TODO Anweisung: ein Satz über der ersten Teilaufgabe fehlt in der Bank – Nr. 18
\begin{aufgabe}{Potenzen – Fehler finden, Begründen}
\swz[3]{Mia rechnet $6^3 = 18$. Wo ist der Fehler?}{}
\swfrage{Begründe, warum $3^5$ größer ist als $5^3$.}
\end{aufgabe}

\uebersichtskasten{Potenz: Eine Potenz ist die Kurzschreibweise für ein Produkt aus lauter gleichen Faktoren – die Basis steht unten, der Exponent (die Hochzahl) sagt, wie oft sie als Faktor vorkommt: 3⁴ = 3 · 3 · 3 · 3 = 81 (nicht 3 · 4 = 12). \\ Exponent finden: Nimm die Basis so oft mit sich selbst mal, bis der Wert erreicht ist, und zähle die Faktoren: $5^x$ = 125 → 5, 25, 125 → x = 3. \\ Negative Basis: Gerade Hochzahl gibt plus, ungerade gibt minus; ohne Klammer gehört das Minus nicht zur Basis: (−2)⁴ = 16, (−2)³ = −8, aber −2⁴ = −16. \\ Negative Hochzahl: Das Minus im Exponenten heißt „eins geteilt durch“ – die Potenz wird ein Bruch, keine negative Zahl: 2⁻³ = 1/2³ = 1/8 = 0,125; jede Basis hoch null ist eins.}
